\section{The horizon, and what is whose}
\label{sec:horizon}

A projection that runs to \ProjectionEndYear{} using inputs that stop at 2100 has
a problem of attribution rather than of arithmetic. Everything past the last year
of published assumptions is the modeller's own extrapolation, and if it is drawn
as one continuous line the reader cannot tell where the source ended. This
section states the boundary explicitly and reports what crosses it.

\subsection{Three separate objects}

\begin{description}[leftmargin=1.4em,style=nextline]
  \item[The UN reproduction.] World Population Prospects 2024's medium path,
  reproduced by this engine, through \OfficialBoundaryYear{} and no further. No
  value after that year is labelled a UN projection.
  \item[The project extension.] Begins from the published population on 1 January
  \OfficialBoundaryYear{} and runs to \ProjectionEndYear{}. Fertility, mortality
  and the sex ratio at birth are held at their final published age schedules,
  which is our assumption and not the source's. Migration continues
  stochastically.
  \item[The selection model.] A separate object that forks in
  \PhaseFiveBaseYear{}, so that selection acts over the whole projection rather
  than being switched on after the UN's assumptions have already determined the
  \OfficialBoundaryYear{} population.
\end{description}

\begin{figure}[t]
\centering
\includegraphics[width=\textwidth]{fig4-boundary-2100.pdf}
\caption{Three objects, drawn so that it is visible which is whose. The solid
grey line is the reproduced UN medium path and stops where the UN's assumptions
stop. The dashed line and shaded fan are this project's own continuation, in
which fertility and mortality schedules are frozen and migration varies across
\ExtensionPaths{} paths; the fan is narrow because migration cancels globally by
construction. The two coloured lines are the selection model, which forks in
\PhaseFiveBaseYear{} and is not a UN projection at any point.}
\label{fig:horizon}
\end{figure}

\subsection{The extension, and why its band is narrow}

The published Bayesian migration trajectories \citep{azose2015} run to 2100 and
the public archive does not include the fitted posterior state, so the
continuation refits the same autoregressive form to the
\ExtensionFitStart{}--\ExtensionFitEnd{} portion of those trajectories and
continues each path from its own \OfficialBoundaryYear{} rate. This is an
emulator of a model's output, not an official continuation, and it is labelled as
one throughout.

The result is \ExtensionTwentyOneHundred{} billion in \OfficialBoundaryYear{},
\ExtensionTwentyOneTwentyFive{} billion in 2125, and
\ExtensionTwentyOneFifty{} billion in \ProjectionEndYear{}, with a
migration-only 90\% range of \ExtensionTwentyOneFiftyLow{} to
\ExtensionTwentyOneFiftyHigh{} billion. The narrowness of that band is a
structural fact rather than a claim of precision: net migration sums to zero
across the world in every draw-year, so it can only move a world total
indirectly, by relocating people into countries whose fertility and mortality
schedules differ from the ones they left.

\subsection{Which uncertainty dominates depends on where you look}

\begin{figure}[t]
\centering
\includegraphics[width=\textwidth]{fig5-decomposition.pdf}
\caption{\textbf{(a)} The 90\% width of world population in \ProjectionEndYear{}
attributable to each source, varying one at a time across its own draws with the
others held at a median trajectory. The post-2100 entry is what our own
hold-constant rule \emph{adds}: it is the extra width from ending the source data
ten years earlier, which is the cleanest available measure of how much of the
answer is ours rather than the source's. \textbf{(b)} The same comparison inside
individual countries at 2100, where the ordering reverses.}
\label{fig:decomposition}
\end{figure}

Varying one source at a time across \DecompositionDraws{} draws
(\Cref{fig:decomposition}a), the 90\% width of world population in
\ProjectionEndYear{} is \WidthFertility{} billion for long-run fertility,
\WidthMechanism{} billion for the selection mechanism's own parameters,
\WidthHoldConstant{} billion for our post-2100 hold-constant rule,
\WidthMortality{} billion for mortality and \WidthMigration{} billion for
migration. The mechanism entry is not directly additive with the others, because
it varies mechanism parameters against the UN's medium demographic rates rather
than against the probabilistic posterior; it is included because a decomposition
that omitted the paper's own subject would be a strange thing to publish.

Two of these deserve comment. Our own post-2100 rule contributes more world
spread than mortality and migration combined, which is a useful corrective: at
this horizon the modeller's housekeeping choices are not negligible relative to
the components everyone argues about. And the selection mechanism is second only
to long-run fertility, which is the case for treating it as a first-class source
of uncertainty rather than an omitted refinement.

The country panel (\Cref{fig:decomposition}b) disagrees with the world panel, and
that disagreement is the point. At 2100, migration uncertainty is
\MigrationOverFertilityUae{} times fertility uncertainty for the United Arab
Emirates and \MigrationOverFertilityCanada{} times for Canada, but only
\MigrationOverFertilityNigeria{} times for Nigeria. In the project extension the
Emirates' \ProjectionEndYear{} population ranges from
\UaeExtensionLow{} million to \UaeExtensionHigh{} million around a median of
\UaeExtensionMedian{} million --- a factor of more than two hundred, for a
component that contributes almost nothing to the world total. Which uncertainty
dominates is a fact about where you are looking, and any single summary of
``the'' uncertainty in a population projection is answering a question about the
world that most users of projections are not asking.

\subsection{A conventional comparator, recorded in advance}
\label{sec:comparator}

For a benchmark that nobody could accuse us of arranging afterwards, the
published Bayesian posterior for fertility and mortality
\citep{alkema2011,raftery2012,sevcikova2016bayespop} was propagated through this
same engine, one draw at a time: \ComparatorDraws{} draws,
\ComparatorCountries{} countries, to \ProjectionEndYear{}. Its median peaks at
\ComparatorPeak{} billion in \ComparatorPeakYear{} and reaches
\ComparatorTwentyOneFifty{} billion in \ProjectionEndYear{}, with a 90\% band of
\ComparatorTwentyOneFiftyLow{} to \ComparatorTwentyOneFiftyHigh{} billion.

This is the conventional mean-reverting view, not this paper's. It was computed,
stored as an immutable dated record, and marked as not a project claim
\emph{before} the mechanism results existed, precisely so that the comparison
could not be assembled to flatter the mechanism afterwards. One check on it is
worth reporting because nobody arranged it: the wholly separate deterministic run
on the UN's own published assumptions peaks at \UnEquivalentPeak{} billion in
\UnEquivalentPeakYear{}, against the ensemble median's \ComparatorPeak{} billion
in \ComparatorPeakYear{}. Those two calculations share only the engine, and they
agree on the peak to within a fifth of a billion people and on its date to within
a decade.

\subsection{What the historical record says about published ranges}

The same archive machinery grades \BacktestVintages{} archived UN revisions from
\BacktestFirstVintage{} to \BacktestLastVintage{} against the current estimates.
World population was under-projected by \BacktestWorldBias{} on average and in
the same direction in every revision since 1996; African fertility was projected
\BacktestAfricaFertility{} too low and East Asian fertility
\BacktestEastAsiaFertility{} too high; life expectancy was
\BacktestLifeExpectancy{} years too low.

We report this not as evidence about our own model, which it is not, and not as a
reason to prefer higher numbers. Its useful content is about ranges rather than
central values: only \BacktestInsideRange{} of \BacktestRangeTotal{}
world-level projections (\BacktestRangeShare{}) fell inside the low-to-high band
printed in the revision that made them. Those bands were fertility scenarios
without probabilities attached, so this is not a calibration failure in the
technical sense. It is a demonstration that the intervals demography has
historically published around long-run population have been too narrow, which is
the case for the kind of explicit uncertainty accounting attempted here.
